\section{Conclusion}
\label{sec:conclusion}

We conducted the first systematic empirical study of linear direction
compositionality in LLM residual streams, testing 10 concept pairs across 5
semantic categories in \pythia. Our results reveal that most mean-difference
concept directions are surprisingly non-orthogonal (median $\abscossim = 0.91$)
and that compositionality depends critically on geometric alignment: near-orthogonal
pairs compose well under steering, while anti-parallel pairs fail
catastrophically. Probing-based composition metrics are misleadingly optimistic,
masking failures that steering exposes. The absolute cosine similarity between
directions is the most reliable predictor of composition success, more
informative than semantic category labels.

The key practical takeaway is simple: \textbf{check orthogonality before
composing steering vectors.} If $\abscossim > 0.5$, naive vector addition will
likely produce lossy or incoherent results. Practitioners should prefer
SAE-based or orthogonalized directions, or apply different vectors at different
layers following \citet{stolfo2024improving}.

Several directions for future work emerge from our findings. First, scaling to
50+ concepts with SAE-based direction extraction~\citep{chalnev2024improving}
would test whether cleaner directions improve composability. Second, using the
causal inner product~\citep{park2023linear} for orthogonality measurement may
reveal structure that standard cosine similarity misses. Third, testing on
larger models~\citep{todd2023function} would establish whether composability
improves with scale. Finally, extending beyond pairwise composition to 3-way
and 4-way combinations would reveal whether composition errors compound, as our
component asymmetry results suggest they might.
